\section*{Chapitre \ref{ch:g11:minerals} --- Les minéraux}

\begin{solution}{exo:g11:minerals:1}
Le \omterm{def:g11:minerals:cleavage}{clivage} est la tendance d'un minéral à se casser selon des plans,
parallèles à des directions dans lesquelles les liaisons entre ses atomes
sont plus faibles~; la \omterm{def:g11:minerals:cleavage}{cassure} est une rupture selon des surfaces
irrégulières, dans un minéral sans plan de faiblesse. Le quartz n'a pas de
\omterm{def:g11:minerals:cleavage}{clivage}~: il montre une \omterm{def:g11:minerals:cleavage}{cassure} conchoïdale, à surfaces lisses et courbes
comme du verre brisé. Un mica a un \omterm{def:g11:minerals:cleavage}{clivage}, parfait~: il se détache en
feuillets.
\end{solution}

\begin{solution}{exo:g11:minerals:2}
Une trace d'un autre élément suffit à changer la couleur du quartz, qui peut
être incolore, blanc, noir, rose, bleu ou violet~; l'olivine est presque
toujours verte. Plus fiables~: le \omterm{def:g11:minerals:cleavage}{clivage} (ou la \omterm{def:g11:minerals:cleavage}{cassure}), la dureté,
l'\omterm{def:g11:minerals:lustre}{éclat}, l'\omterm{def:g11:minerals:habit}{habitus} des cristaux libres, la masse volumique.
\end{solution}

\begin{solution}{exo:g11:minerals:3}
(1) Olivine. (2) Biotite. (3) Orthose (un feldspath~: deux \omterm{def:g11:minerals:cleavage}{clivages}, dureté
6). (4) Quartz (pas de \omterm{def:g11:minerals:cleavage}{clivage}, dureté 7).
\end{solution}

\begin{solution}{exo:g11:minerals:4}
(1) Cinq atomes~: un silicium au centre et quatre oxygènes aux sommets d'un
tétraèdre. (2) $+4 + 4 \times (-2) = -4$~: le groupe est \ce{SiO4^4-}.
\end{solution}

\begin{solution}{exo:g11:minerals:5}
Olivine~: tétraèdres isolés. Pyroxènes~: chaînes simples. Amphiboles~:
chaînes doubles. Micas~: feuillets. Feldspaths et quartz~: charpentes.
\end{solution}

\begin{solution}{exo:g11:minerals:6}
(1) $M(\ce{SiO2}) = 28.085 + 2 \times 15.999 = \qty{60.083}{g/mol}$~;
oxygène~: $31.998/60.083 \approx \qty{53.3}{\percent}$.
(2) $M(\ce{Mg2SiO4}) = 2 \times 24.305 + 28.085 + 4 \times 15.999 =
\qty{140.691}{g/mol}$~; oxygène~: $63.996/140.691 \approx \qty{45.5}{\percent}$.
(3) La croûte, à \qty{46}{\percent}, est du même ordre~: elle est faite
presque entièrement de minéraux riches en oxygène, les \omterm{def:g11:minerals:silicate}{silicates}.
\end{solution}

\begin{solution}{exo:g11:minerals:7}
Le grain blanc a deux \omterm{def:g11:minerals:cleavage}{clivages} et une dureté comprise entre 5.5 et 6.5~:
c'est un feldspath (6). Le grain gris est plus dur que 6.5 et montre une
\omterm{def:g11:minerals:cleavage}{cassure} courbe, conchoïdale, sans \omterm{def:g11:minerals:cleavage}{clivage}~: c'est du quartz (7).
\end{solution}

\begin{solution}{exo:g11:minerals:8}
(1) $4 - 2/2 = 3$ oxygènes par silicium~: \ce{SiO3}. (2) \ce{(Si,Al)2O6}
contient deux atomes de silicium (ou d'aluminium) pour six oxygènes, trois
par silicium~: c'est deux fois \ce{SiO3}.
\end{solution}

\begin{solution}{exo:g11:minerals:9}
$\rho = 0.55 \times 2.69 + 0.35 \times 3.4 + 0.10 \times 3.32 = 1.480 + 1.190
+ 0.332 \approx \qty{3.00}{g/cm^3}$. Le cube a un volume de
\qty{1000}{cm^3}, donc une masse d'environ \qty{3.00}{kg}.
\end{solution}

\begin{solution}{exo:g11:minerals:10}
(1) $0.75 \times 3.32 + 0.25 \times 3.4 = 2.49 + 0.85 = \qty{3.34}{g/cm^3}$.
(2) $(3.34 - 3.00)/3.00 \approx 0.11$~: environ \qty{11}{\percent} plus
dense. (3) Le basalte et le \omterm{def:g11:ocean-ridges:gabbro}{gabbro} de la \omterm{def:g11:ocean-ridges:oceanic-lithosphere}{croûte océanique} sont plus légers
que la \omterm{def:g11:ocean-ridges:peridotite}{péridotite} du manteau~: ils restent au-dessus d'elle.
\end{solution}

\begin{solution}{exo:g11:minerals:11}
(1) Plagioclase $168/400 = 0.42$, pyroxène $0.33$, olivine $0.15$, trous
$0.10$. (2) La proposition vaut pour une roche sans pores~; ici, un dixième
du volume est vide. (3) Dans la partie solide, les fractions valent
$0.42/0.90
\approx 0.467$, $0.367$ et $0.167$~:
$\rho_{\text{solide}} \approx 0.467 \times 2.69 + 0.367 \times 3.4 + 0.167
\times 3.32 \approx \qty{3.06}{g/cm^3}$. Avec ses trous, qui ajoutent du
volume et pas de masse, la roche a pour masse volumique
$0.42 \times 2.69 + 0.33 \times 3.4 +
0.15 \times 3.32 \approx \qty{2.75}{g/cm^3}$, c'est-à-dire $0.9 \times 3.06$.
\end{solution}

\begin{solution}{exo:g11:minerals:12}
(1) Oxygène $464/16.0 = 29.0$~mol~; silicium $283/28.1 \approx 10.1$~;
aluminium $84/27.0 \approx 3.1$~; fer $52/55.8 \approx 0.93$~; calcium
$46/40.1 \approx
1.15$~; magnésium $28/24.3 \approx 1.15$~; sodium $23/23.0 = 1.0$~;
potassium $15/39.1 \approx 0.38$. (2) En tout, environ $46.8$~mol
d'atomes, dont
$29.0/46.8 \approx 0.62$~: presque deux atomes sur trois sont des atomes
d'oxygène. (3) Les gros ions oxygène forment la plupart des atomes, et une
part plus grande encore du volume~: à l'échelle des atomes, la croûte est un
empilement d'ions oxygène, les petits ions silicium et métalliques occupant
les espaces qui les séparent.
\end{solution}

\begin{solution}{exo:g11:minerals:13}
(1) Le nombre moyen de sommets partagés vaut $(2 + 3)/2 = 2.5$~: il y a donc
$4 - 2.5/2 = 2.75$ oxygènes par silicium. (2) Quatre siliciums pour onze
oxygènes~: \ce{Si4O11}. (3) \ce{Si7AlO22} contient huit tétraèdres (sept de
silicium, un d'aluminium) pour 22 oxygènes, soit $22/8 = 2.75$ par
tétraèdre, ou deux fois \ce{Si4O11}. Les deux groupes \ce{OH} ne font pas
partie des tétraèdres.
\end{solution}

\begin{solution}{exo:g11:minerals:14}
$0.20 \times 2.62 + 0.10 \times 2.56 + 0.50 \times 2.65 + 0.10 \times 3.09 +
0.10 \times 3.23 = 0.524 + 0.256 + 1.325 + 0.309 + 0.323 \approx
\qty{2.74}{g/cm^3}$, au-dessus de \qty{2.7}{g/cm^3}~; le granite de
l'exemple, plus riche en quartz et en feldspath potassique, est à environ
\qty{2.59}{g/cm^3}, sous \qty{2.6}{g/cm^3}. Les roches de l'ouest sont plus
denses parce qu'elles contiennent plus de plagioclase et davantage de
biotite et de hornblende, riches en fer et en magnésium~; les granites
riches en silice contiennent plus de quartz et de feldspath potassique, les
plus légers des \omterm{def:g11:minerals:rock-forming}{silicates constitutifs des roches}.
\end{solution}

\begin{solution}{exo:g11:minerals:15}
(1) Croûte supérieure~: $66.6 \times 0.467 \approx \qty{31.1}{\percent}$ de
silicium~; croûte entière~: $60.6 \times 0.467 \approx \qty{28.3}{\percent}$.
(2) La croûte supérieure, plus riche en silice, devrait contenir plus de
quartz et de feldspaths, et moins de minéraux riches en fer et en
magnésium~; la croûte plus profonde, l'inverse. (3) Le quartz et les
feldspaths sont les plus légers des \omterm{def:g11:minerals:rock-forming}{silicates constitutifs des roches}~: la
croûte supérieure devrait être plus légère que la croûte plus profonde.
\end{solution}

\begin{solution}{pb:g11:minerals:1}
\textbf{1.} Quartz~: pas de \omterm{def:g11:minerals:cleavage}{clivage}, \omterm{def:g11:minerals:cleavage}{cassure} conchoïdale, \omterm{def:g11:minerals:lustre}{éclat} vitreux,
plus dur que le clou (7).

\textbf{2.} Orthose, un feldspath potassique~: deux \omterm{def:g11:minerals:cleavage}{clivages}, dureté 6 sur
l'échelle de Mohs.

\textbf{3.} Un plagioclase (ici l'oligoclase), le feldspath à sodium
et à calcium.

\textbf{4.} D est de la biotite, E de la muscovite~: deux micas. Leurs
tétraèdres forment des feuillets, et les liaisons entre les feuillets sont
faibles~: ils se fendent selon cet unique \omterm{def:g11:minerals:cleavage}{clivage} parfait.

\textbf{5.} Le quartz et les deux feldspaths (A, B, C) sont des
charpentes~; la biotite (D) contient du fer et du magnésium.

\textbf{6.} Quartz $150/500 = 0.30$~; orthose $0.35$~; oligoclase $0.25$~;
biotite $0.06$~; muscovite $0.04$.

\textbf{7.} Quand les grains sont répartis sans ordre, la part des points
qui tombent sur un minéral estime la part de la surface qu'il couvre sur la
face, et celle-ci est égale à la part du volume qu'il occupe.

\textbf{8.} $0.35 + 0.25 = 0.60$~: \qty{60}{\percent} de feldspath.

\textbf{9.} De $1/500 = 0.002$, soit \qty{0.2}{\percent}.

\textbf{10.} Oui~: feldspath, quartz et mica.

\textbf{11.} $\rho = \phi_1 \rho_1 + \phi_2 \rho_2 + \dots + \phi_k \rho_k$,
avec $\phi_1 + \dots + \phi_k = 1$.

\textbf{12.} Quartz $0.30 \times 2.62 = 0.786$~; orthose $0.35 \times 2.56 =
0.896$~; oligoclase $0.25 \times 2.65 = 0.6625$~; biotite $0.06 \times 3.09
\approx 0.185$~; muscovite $0.04 \times 2.82 \approx 0.113$ (tous en
\unit{g/cm^3}).

\textbf{13.} $0.786 + 0.896 + 0.6625 + 0.185 + 0.113 \approx
\qty{2.64}{g/cm^3}$.

\textbf{14.} L'orthose, avec $0.896$, parce que c'est le minéral le plus
abondant. La biotite et la muscovite, à \qty{3.09}{} et \qty{2.82}{g/cm^3},
sont plus denses que les autres minéraux et portent la masse volumique de
la roche au-dessus de celle du quartz et du feldspath seuls.

\textbf{15.} \qty{2640}{kg/m^3}~: \qty{1}{m^3} de ce granite a une masse
d'environ \qty{2.64}{t}.

\textbf{16.} $100 \times 2.643 \approx \qty{264.3}{g}$.

\textbf{17.} Le bloc déplace \qty{100}{cm^3} d'eau, de masse
\qty{100}{g}~: la balance indique $264.3 - 100 = \qty{164.3}{g}$.

\textbf{18.} $264.3/(264.3 - 164.3) = 264.3/100 \approx \qty{2.64}{g/cm^3}$.

\textbf{19.} Entre les deux~: plus léger que la moyenne de
\qty{2.69}{g/cm^3}, au-dessus des \qty{2.6}{g/cm^3} des granites les plus
riches en silice et en potassium. Davantage de plagioclase et davantage de
biotite ou de hornblende, et moins de quartz et de feldspath potassique,
rendraient un granite plus dense.

\textbf{20.} La masse volumique du granite vaut environ
\textbf{\qty{2.64}{g/cm^3}}.
\end{solution}
